% ─────────────────────────────────────────────────────────────────
%  Curriculum Vitae — \UserName{}
%  Font     : Bitstream Charter (BT Roman)
%  ATS      : Single-column, paracol skills, plain headings
%  Icons    : fontawesome5
%  Footer   : Last Modified date — light gray, bottom right
% ─────────────────────────────────────────────────────────────────
\documentclass[10pt, a4paper]{article}

%% ── Packages ─────────────────────────────────────────────────────
\usepackage[T1]{fontenc}
\usepackage[utf8]{inputenc}
\usepackage{charter}
\usepackage{microtype}
%\usepackage[margin=1.4cm, top=1.4cm, bottom=1.8cm]{geometry}
\usepackage[hmargin=1.25cm, vmargin=0.75cm]{geometry}
\usepackage{parskip}
\usepackage{titlesec}
\usepackage{enumitem}
\usepackage{xcolor}
\usepackage{tabularx}
\usepackage{paracol}
\usepackage{fontawesome5}
\usepackage{fancyhdr}
\usepackage{hyperref}

%% ── Colours ──────────────────────────────────────────────────────
\definecolor{accentblue}{HTML}{2B6CB0}
\definecolor{lightgray}{HTML}{555555}
\definecolor{datecolor}{HTML}{CCCCCC}
\definecolor{iconcolor}{HTML}{2D3748}
\definecolor{ltgrey}{gray}{0.50}

%%%
%%%
% Load Confidential data
\IfFileExists{./data/UserPvtInfo1.tex}{\input{./data/UserPvtInfo1.tex}}{\input{./data/UserPubInfo.tex}}

%% ── Hyperlinks ───────────────────────────────────────────────────
\hypersetup{
  colorlinks = true,
  urlcolor   = accentblue,
  linkcolor  = accentblue,
  pdfauthor  = {\UserName{}},
  pdftitle   = {\UserName{} - Embedded Systems CV},
  pdfsubject = {Curriculum Vitae}
}

%% ── Section style ────────────────────────────────────────────────
\titleformat{\section}
  {\large\bfseries\color{accentblue}}
  {}{0em}{}
  [\vspace{-0.5em}\rule{\linewidth}{0.6pt}\vspace{-0.3em}]
\titlespacing{\section}{0pt}{8pt}{4pt}

\pagestyle{empty}

%% ── List settings ────────────────────────────────────────────────
\setlist[itemize]{leftmargin=1.5em, topsep=2pt, itemsep=1pt, parsep=0pt}

%% ── CV Entry command ─────────────────────────────────────────────
\newcommand{\cventry}[3]{%
  \noindent
  \begin{tabularx}{\linewidth}{@{}X r@{}}
    \textbf{#1} --- \textit{#2} & \textcolor{lightgray}{\small #3}
  \end{tabularx}
  \vspace{2pt}
}

\newcommand{\clientinfo}[2]{%
  \noindent\textcolor{lightgray}{\small\textit{#1:}}\enspace#2\par
}

\newcommand{\MainHeading}[1]{%
    {\noindent \color{ltgrey}\normalsize{\textsc {#1}}}\vspace{2pt}}%%

\newcommand{\MainSubHeading}[2]{%
   {\noindent \large{\textsc{\textbf{#1} #2}}}\\ }

%% ── Bullet Point ─────────────────────────────────────────────
\newcommand{\bpoint}[1]{\mbox{{\textbullet{} #1 }}}


%%%%%%%%%%%%%%%%%%%%%%
%
%
%
\newcommand{\BAspaceBullet}{\quad\textbar\quad}
\usepackage{layouts}
%%%%%%%%%%%%%%%%%%%%%%%%%%%%%%%%%
%% ─────────────────────────────────────────────────────────────────
\begin{document}
%% ══════════════════════════════════════════════════════════════════
%%  HEADER
%% ══════════════════════════════════════════════════════════════════
\begin{minipage}[b]{0.66\textwidth}        % Minipage align to bottom
%\raggedright{
   \centering{
         \UserName{} \\
         Embedded Architect \textbullet{} Linux Kernel \textbullet{} XBL \textbullet{} U-Boot \textbullet{} UEFI \textbullet{} UFS \newline %%
         Ethernet \textbullet{} KPI optimizations \newline %%
         Automotive ADAS/IVI \vspace{2pt}\newline %%
         \small{\faPhone*{} \UserMobile{} \quad\textbar\quad \faEnvelope{} \space \UserEmailID{} \quad\textbar\quad \faLinkedin{} \space \UserLinkedIn{}} %%
   } %% \centering
\end{minipage}%
\hfill
%
%   Begin Left side Minipage
   %\noindent\fbox{%
\begin{minipage}[b]{0.33\textwidth}        % Minipage align to bottom
      \raggedleft{
         \CurrentCompany{} \\%%
            Staff Engineer  \\%%
            Exp. \UserExperience{} Years\\%%
            \UserLocation{}\\%%
      }%%
\end{minipage}%
%\vspace{2pt}\rule{\linewidth}{1.2pt}\vspace{4pt}
%%
%% ══════════════════════════════════════════════════════════════════
%%  PROFESSIONAL SUMMARY
%% ══════════════════════════════════════════════════════════════════
\vspace{1em}%%
\newline%%
\MainHeading{Professional Summary}\vspace{-5pt}%%
\begin{itemize}[label={\textbullet}, leftmargin=*, topsep=0pt, partopsep=0pt, itemsep=0pt, parsep=0pt] %%
   \item Nearly 13 years of experience in Embedded Software Development, contributing to the design, development, debugging, and support of software products across the complete development lifecycle.
   \item Strong technical background in Linux kernel, bootloaders, schematic, storage technologies, embedded protocols, device drivers, and system-level software development.
   \item Driving software architecture and design, technical discussions, scoping, ETAs and tracking execution.
   \item Worked closely with customers, product management, program management, and cross-functional teams to define requirements, execution plan, resolve issues, and deliver solutions.
   \item Recognized for hands-on problem solving, root cause analysis, performance optimization, and a collaborative approach to improving product quality and reliability.
\end{itemize}
%
%
%
%
%
%% ══════════════════════════════════════════════════════════════════
%%  TECHNICAL SKILLS
%% ══════════════════════════════════════════════════════════════════
\vspace{1em}%%
\begin{minipage}[t]{0.33\textwidth}        % Minipage align to bottom
%% =========== Point wise skills =============
\raggedright{
   \MainHeading{Core Competencies}\newline%%
   {
      \bpoint{Software Architect}
      \bpoint{Linux kernel}
      \bpoint{XBL}
      \bpoint{U-Boot}
      \bpoint{UEFI}
      \bpoint{Trace32}
      \bpoint{HW Schematic}
      \bpoint{UFS}
      \bpoint{I2C}
      \bpoint{PIN Multiplexing}
      \bpoint{SPI}
      \bpoint{MDIO}
      \bpoint{Board Bring up}
      \bpoint{Ethernet}
      \bpoint{Switch - MAC Layer 2 \& Layer 4}
      \bpoint{Board Bring-UP Activities}
      \bpoint{Boot Flow - PBL to HLOS}
      \bpoint{System Development}
      \bpoint{IPC}
      \bpoint{Yocto}
      \bpoint{git} \bpoint{perforce} \bpoint{svn} \bpoint{Jenkins} \bpoint{Trace32}
      \bpoint{Project Management}
      \bpoint{Team Lead}
      \bpoint{Claude workflows}
   }
} %%\raggedright
\newline %%

%% ══════════════════════════════════════════════════════════════════
%%  EDUCATION
%% ══════════════════════════════════════════════════════════════════
   \MainHeading{Education}\newline%%
   \MainSubHeading{B.Tech}{in Electronics \& Communication Engg.}
   {\textsc {\small \color{ltgrey}\UserCollege{} \hfill Apr 2013}}
\newline %%

   \MainHeading{Programming}\newline%%
   \bpoint{Proficient in C \& Data Structures}
   \bpoint{Bash}
   \bpoint{Python}
   \bpoint{C++}
\newline %%


   \MainHeading{Languages}\newline%%
   \bpoint{Native \emph{Panjabi}}
   \bpoint{Proficient in English}
   \bpoint{Hindi}
\end{minipage}%
%
\hfill
%
\begin{minipage}[t]{0.66\textwidth}        % Minipage align to bottom
  \MainHeading{Technical Skills}%%
   \begin{itemize}[label={-}, topsep=0pt, partopsep=0pt, itemsep=0pt, parsep=0pt]
%
% ============== Leadership =============== %
      \item \noindent Implemented Agentic AI workflow with Claude, created workflows for storage.
      \item \noindent Expertise in end-to-end development of software products from requirement analysis to system study,
       designing, coding, testing, debugging, documentation and implementation.
      \item \noindent Optimizing the existing design, software, Improving Boot KPIs i.e. boot/driver to add value to product.
      \item \noindent Good Analytical skills to segregate complex issues for systematic resolution.
      \item \noindent Supporting customers for blocker and escalated issues providing timely support and unblocking them.
      \item \noindent Leading failure triage, root cause analysis, and system level debugging across platforms.
      \item \noindent Driving project planning, task breakdown, estimation, prioritization, and execution tracking.
      \item \noindent Complete understanding of the boot flow and involvement of different stages involved till final application.
%
% ============== Technical =============== %
      \item \noindent Expertize in embedded protocols (UFS, I2C, SPI, MDIO, Ethernet L2) and Device driver development in Linux Kernel and Boot loader i.e. u-boot, XBL, UEFI.
%      \item \noindent Working on Linux kernel and Zephyr OS.
      \item \noindent Experience in MAC driver for ethernet at L2.
      \item \noindent Knowledge of Network Pakcet flow and with TCP/IP layer.
%      \item \noindent Planned the software architecture as per customer's requirement.
      \item \noindent Done the Board Bring up, porting of u-boot, device-tree and linux.
      \item \noindent Experience in debugging of interface like DDR, NAND, I2C, SPI and UART.
      \item \noindent Experience in networking for L2, priority scheduling, Round Robin scheduling.
      \item \noindent Experience in TCP/IP, UDP, layer 2, layer 4 and complete network packet flow.
      \item \noindent Experience in Ethernet driver, socket buffers i.e. skb packets to prioritize UDP packets.
      \item \noindent Write slave drivers for I2C and SPI for micro-processor, as Linux kernel doesn’t provide support for slave mode.
      \item \noindent Write linux device driver to access DDR for firmware upgrade \& fail-safe mechanism.
      \item \noindent Experience with interfaces like NAND, EMMC, Ethernet, USB, DDR, NOR, Sd-Card, Bluetooth.
      \item \noindent Experience in debugging I2C, SPI, UART slave devices and at linux user space.
      \item \noindent Customized boot-loader and Linux operating system for embedded platforms.
      \item \noindent Experienced with i.MX 6 sabresd/sololite, Mini2440,
         \href{http://www.radiumboards.com/TI\_TDA2x\_Based\_xCAM\_Platform.php}{\textbf{\underline{TI's TDA2x}}},
         \href{http://www.nxp.com/products/reference-designs/qoriq-ls1021a-iot-gateway-reference-design:LS1021A-IoT}{\textbf{\underline{LS1021A-IOT}}} boards.
      \item \noindent Experience with processors like TI's omap, Freescale's i.MX 6, Qualcomm's IPQ4029 \& IPQ4029.
      \item \noindent Ported the Marvell 88E6352 chip in linux and u-boot based on Distributed Switch Architecture for Ethernet.
      %\item \noindent Experience with SDK: Yocto, Openwrt, VISION\_SDK.
      %\item \noindent Having a good analysis and debugging skills.
   \end{itemize}
\end{minipage}%
%
%
%
%
%
%% ══════════════════════════════════════════════════════════════════
%%  PROFESSIONAL EXPERIENCE
%% ══════════════════════════════════════════════════════════════════
\newline%%
\vspace{0pt}%%
\MainHeading{Professional Experience}\newline%%
%% ── Job 5 ────────────────────────────────────────────────────────
\cventry{Staff Engineer}{\CompanyV{}}{\CompanyVExp{}}
%\clientinfo{Location}{\CompanyVLoc{}}
\vspace{-1.5em}%%
\begin{itemize}[label={-}, leftmargin=*, topsep=0pt, partopsep=0pt, itemsep=0pt, parsep=0pt] %%
   \item Overall Storage technology Lead for all the AUTO SOC’s primarily for UFS \& its dependency across subsystem like PBL, PMIC, IMEM, Hypervisor, TZ, DDR and boot.
   \item Involved in end-to-end design from requirement analysis to system study, limitations, active discussions across the teams.
   \item Handling Feature requests, development, triaging blocker and customer issues.
%Expertise in end-to-end development of software products from requirement analysis to system study, designing, coding, testing, debugging, documentation and implementation.
   \item Planning the project, task breakdown, discussion across the teams, dependencies, prioritization and execution tracking
   \item Effectively addressing stakeholders across management, CE, Customers and internal teams for the timelines.
\end{itemize}

\begin{itemize}
  \item \textbf{WakeUP sequence broken for Suspend to Ram feature}
		Resolved crash in UFS linux kernel while waking up the system during Suspend to RAM
		feature. Analyzed the state machine and different subsystems i.e. Linux, GearVM
		involved in this use-case to figure out the crash due to power change event and system
		remains in hibernation state. Changes mainlined in gearvm for scalability and upstream changes.

  \item \textbf{90msec improvement by skipping UFS init in kernel}
		Designed and implemented feature to skip UFS init in linux kernel to leverage the previous
		state machine as configured through XBL, UEFI. At XBL, UFS is initialized to HS-G4 and maintined
		till linux kernel. Make sure the UFS is not put to SSU or in reset state. Handled UEFI Exit
		Boot Services to avoid SSU trigger for UFS.

  \item \textbf{Boot KPI Optimisation --- 200+ msec Saved:}
		Optimized existing boot flow that resulted in 200+ msec improvement
		• UFS partial Init • Remove retry of UFS2.x
		• Remove UFS read in PWM Gear for GPT • UFS read size increased to 256MB data
		• PMIC init parallel to UFS

  \item \textbf{Scatter Gather --- 70\% KPI Improvement:}
        Designed and implemented \textbf{Scatter Gather} for UFS read operations,
        enabling a single read command to load boot images to multiple discontiguous
        DDR addresses, replacing multiple sequential reads and achieving a
        \textbf{70\% improvement} in boot KPI.

  \item \textbf{2nd UFS initialization:}
        Initialize the 2nd UFS in UEFI, device programmer and linux kernel. Design take care of
		KPI impact, Hyervisor passthrouh access, customer scenarios as 2nd UFS is not PORd for all customers. Analyze the
		pre-requisities like PON sequence, Clocks and voltages.

  \item \textbf{UFS PHY Skip Feature:}
        Designed the UFS subsystem state to make PHY configurations redundant at XBL to save \textbf{100 msec}.
		PBL will configure the PHY registers with HS Gear settings and 2 lanes but still operate over
		1 Lane that makes PHY configurations redundant at XBL which requires only Gear Change and
		Data lanes to 2. Proof of Concept done at pre-sil stage and identified the critical crash at XBL as 
		\texttt{PA\_AvailLanes} was not set correctly in PBL.


  \item \textbf{Mission \& Non-Mission Mode (Licensed Feature):}
        Designed and delivered a \textbf{licensed JEDEC-spec feature} for subscribed
        AUTO OEMs enabling comprehensive UFS device validation in UEFI, providing
		access to JEDEC defined descriptor, attributes and flags. It helps to analyze
		state machine in case of crash observed for the on field devices.

  \item \textbf{UFS Cache Bug --- 50K Iteration Failure:}
        Customer Issue - Diagnosed and fixed a critical intermittent failure (per 50K iterations)
        across multiple UFS commands (bootable LUN, geometry, config, unit
        descriptors) caused by \textbf{stale cache responses} on the Host
        Controller; flushing cache memory resolved all failure scenarios. Analyze the
		pattern across failures due to mismatch in the tag of sent request 
		and received response. Provided patch to customer quickly.

  \item \textbf{PMIC controller to set VCCQ to 1.224 --- Critical Customer issue:}
		Customer devices requires VCCQ to 1.224V than 1.2V from reference HW and only when UFS3.x
		card than UFS2.x. Impact analysis and design shared with customer on XBL and Flashing sequence.
		In XBL, as UFS is already initialized the UFS so Spec Version can be read to init the PMIC first and then full PHY init
		happens. In Device Programmer, It requires re-int

  \item \textbf{First UFS Read Failure After PMIC Init:}
        Investigated a critical bug where the first UFS read command failed
        after PMIC initialisation with PA/DL layer errors. Traced root cause to
        \textbf{differing PSI voltage values} between PON sequence and
        post-PMIC-init state causing voltage jitter; routed confirmed issue to
        PMIC team.

  \item \textbf{Automotive SoC Bring-Up --- SA8255 (LeMans):}
        Led the UFS bring-up for \textbf{SA8255 (Qualcomm's first AUTO SoC)},
        meticulously planning all XBL, UEFI, and flashing tool use-cases,
        dependencies, and limitations. Resolved critical clock issues while
		configuring PHY registers, provided work around for different issues
		(Scale PBL settings with 1 Lane, enable all clocks)to unblock further
		development and parallely working with relevant teams for proper fix.
		Also successfully brought up UFS on \textbf{SA8775, SA8620, and SA8797}.

  \item \textbf{BU - SA8797, NOP failure:}
        Resolved the SBFE i.e. System Bus Fatal Error reported at UFS Host upon NOP command
		is being ready to sent. UFS Host controller facing issues while access the memory
		where command is being prepared. Debugged and provided workaround to allocated memory
		on OCIMEM rather than usual IMEM. Further, issue is resolved as memory is not 4KB aligned.

  \item \textbf{GearVM: Architect GearVM Software responsibilities}
		Architected and implemented the UFS subsystem in GearVM to manage performance
		and power-domain requests from QNX, GVM, and Linux clients. Designed a granular
		power-management framework enabling OS to independently control	UFS states,
		including Power On, Power Off, Sleep, and Suspend-to-RAM. This architecture
		abstracts Qualcomm-specific and hardware-dependent functionality within 
		GearVM, improving Linux upstream compatibility and reducing platform-specific dependencies.
		Defined and communicated a phased implementation roadmap and progress 
		updates to stakeholders across the program lifecycle.

  \item \textbf{Production Line issue GearVM: Architect GearVM Software responsibilities}
		Customer UFS devices are pre-provisioned at factory with wrong information that enables
		only LUN 2,3 and 7 that cause failure in flashing which furhter blocks the production line.
		Shared analysis with customer that rules out Qualcomm.

  \item \textbf{Cross-Team Leadership:}
        Leading design and development for UFS across interdependent teams
        (PMIC, Clocks, PBL, XBL, Boot); supporting CE teams globally across
        multiple time zones for storage-related customer escalations.

  \item \textbf{Team Leadership:}
        Leading a team of 4 members that covers task planning, mentoring, limitations,
		helping them in debug, navigating there tasks and ETAs.

\end{itemize}

\vspace{4pt}

%% ── Job 4 ────────────────────────────────────────────────────────
\cventry{Technical Leader}{\CompanyIV{}}{\CompanyIVExp{}}
%\clientinfo{Location}{\CompanyIVLoc{}}
\clientinfo{Domain}{AUTOSAR SWS Specification, RADAR Validation, Microcontroller Bring-Up}

\begin{itemize}
  \item \textbf{AUTOSAR SWS Specification:}
        Defined \textbf{AUTOSAR SWS specifications} for Port, SPI, I2C, and
        peripheral interfaces for a new microcontroller, ensuring critical
        timing requirements, safety margins, and implementation guidelines.

  \item \textbf{CAN 50 msec POR Requirement:}
        Analysed and implemented the requirement to initiate the
        \textbf{first CAN message within 50 msec} after Power-On Reset,
        meeting stringent automotive real-time constraints.

  \item \textbf{RADAR Chipset Integration (76--81 GHz):}
        Integrated a RADAR chipset over I2C bus with the \textbf{NXP S32R} series
        microcontroller; implemented full test-case coverage across the complete
        76--81 GHz frequency band with configurable sweep windows.

  \item \textbf{RADAR Matlab Plot Analysis:}
        Analysed generated Matlab plots to validate RADAR test coverage and
        confirm compliance with automotive safety frequency requirements.
\end{itemize}

\vspace{4pt}

%% ── Job 3 ────────────────────────────────────────────────────────
\cventry{Technical Leader}{\CompanyIII{}}{\CompanyIIIExp{}}
%\clientinfo{Location}{\CompanyIIILoc{}}
\clientinfo{Domain}{5G mmWave Wireless, OTA Firmware Upgrade, Ethernet/PHY Integration}

\begin{itemize}
  \item \textbf{Fail-Safe OTA Firmware Upgrade:}
        Implemented a \textbf{fail-safe recovery mechanism} for OTA firmware
        updates targeting boot binaries; ensured system integrity is maintained
        even upon partial or failed OTA operations.

  \item \textbf{Atheros PHY Integration:}
        Ported and integrated \textbf{Atheros PHY} with Marvell SoC using MDIO
        protocol in Linux kernel; provided full software support to meet
        passing criteria.

  \item \textbf{DFS RADAR Detection --- Wireless Client (First-Ever Design):}
        Designed and implemented a wireless client use-case to abort 5\,GHz
        communication and raise CAC to station upon RADAR detection on the
        current DFS channel --- the \textbf{first-ever design} of this type for
        a wireless client, required for Airport safety approvals.
\end{itemize}

\vspace{4pt}

%% ── Job 2 ────────────────────────────────────────────────────────
\cventry{Senior Engineer}{\CompanyII{}}{\CompanyIIExp{}}
%\clientinfo{Location}{\CompnayIILoc{}}
\clientinfo{Domain}{Boot Automation, Jenkins, Wi-Fi Firmware Loading}

\begin{itemize}
  \item \textbf{Automated Boot-Terminal Firmware Loader:}
        u-boot: Implemented a boot terminal command to \textbf{auto-initialise localhost
        IP, fetch Wi-Fi firmware binary over TFTP}, and flash it to the required
        memory address --- fully automating firmware loading at boot time.

  \item \textbf{Jenkins Framework Scaling:}
        Studied and extended the existing Jenkins CI framework to support
        \textbf{new hardware targets} for nightly regression builds and tests.
\end{itemize}

\vspace{4pt}

%% ── Job 5 ────────────────────────────────────────────────────────
\cventry{Senior Engineer}{\CompanyI{}}{\CompanyIExp{}}
%\clientinfo{Location}{\CompanyILoc{}}
\clientinfo{Domain}{Embedded Systems, Board Bring-Up, Linux Kernel, Networking}

\begin{itemize}
  \item \textbf{Multi-Target Board Bring-Up:}
        Led board bring-up for \textbf{TI DM388, Qualcomm IPQ4029, TI TDA2xx,}
        and \textbf{Freescale i.MX6} series; performed U-Boot porting,
        device-tree customisation, and Linux kernel optimisations.

  \item \textbf{Marvell 88E6352 Switch Porting:}
        Ported the \textbf{Marvell 88E6352} Gigabit Ethernet switch chip in both
        Linux and U-Boot based on the \textbf{Distributed Switch Architecture (DSA)},
        enabling multi-port Ethernet switching for IoT router platforms.

  \item \textbf{I2C / SPI Slave Drivers:}
        Written Linux slave drivers for I2C and SPI on microprocessors where
        the Linux kernel does not natively provide slave-mode support.

  \item \textbf{DDR Access Driver --- OTA \& Fail-Safe:}
        Developed a Linux device driver for direct DDR access to support
        firmware upgrade and fail-safe recovery mechanisms.

  \item \textbf{Ethernet MAC Driver (L2):}
        Developed and debugged Ethernet MAC driver at Layer 2, including
        socket buffer (SKB) management for UDP packet prioritisation,
        L2 priority scheduling, and Round Robin scheduling.

  \item \textbf{System Applications:}
        Written applications for \textbf{hostapd, udhcpc, wpa\_supplicant,
        Node.js, and BOA server} to meet end-customer product use-cases;
        implemented IPC, multi-thread, and socket programming solutions.

  \item \textbf{Full-Cycle Ownership:}
        Owned complete embedded system lifecycle: customer requirement
        analysis, feature development, bug resolution, and end-to-end
        validation across bring-up, U-Boot, Linux kernel, and user-space.
\end{itemize}

%% ══════════════════════════════════════════════════════════════════
%%  KEY ACHIEVEMENTS
%% ══════════════════════════════════════════════════════════════════
\MainHeading{\color{accentblue}Key Achievements}\vspace{-1.5em}\newline %%
\begin{itemize}
  \item \textbf{70\% Boot KPI Improvement via Scatter Gather:}
        Designed and delivered Scatter Gather for UFS read, replacing multiple
        sequential reads with a single operation and achieving a
        \textbf{70\% improvement} in boot KPI on Qualcomm AUTO SoCs.

  \item \textbf{200+ msec Cumulative Boot Optimisation at Qualcomm:}
        Drove a portfolio of UFS boot optimisation initiatives (PHY Skip, Probe
        Optimisation, PMIC Parallelisation, PWM Gear removal) delivering
        \textbf{200+ msec total boot KPI reduction} for Automotive SoCs.

  \item \textbf{SA8255 (LeMans) --- First-Ever AUTO SoC UFS Bring-Up:}
        Successfully led UFS bring-up for Qualcomm's first Automotive SoC
        (SA8255/LeMans) and subsequently for SA8775, SA8620, and SA8797.

  \item \textbf{First-Ever DFS RADAR Detection Design for Wireless Client (PHAZR):}
        Designed a novel wireless client use-case for DFS RADAR detection
        in the 5\,GHz band, the first of its kind for a wireless client platform,
        required for airport safety compliance.

  \item \textbf{Licensed Mission Mode Feature for AUTO OEMs:}
        Architected and delivered a JEDEC-spec licensed UFS validation feature
        for subscribed Automotive OEMs enabling comprehensive UEFI-level
        device health testing.
\end{itemize}

%% ── ATS keyword block (white text — invisible to human reader) ───
\begin{center}
\textcolor{white}{\tiny
Linux Kernel, Embedded Linux, Device Driver, U-Boot, UEFI, Bootloader,
UFS, eMMC, NAND, NOR, Storage, JEDEC, JESD220, PHY, HCI, UTRD, UPIU,
Scatter Gather, Write Booster, Provisioning, Automotive, ADAS, SoC,
SA8255, SA8775, SA8620, SA8797, Qualcomm, Staff Engineer, Software Architect,
AUTOSAR, SWS, CAN, RADAR, 76GHz, 81GHz, DFS, 5GHz, mmWave,
Ethernet, MAC, TCP/IP, UDP, L2, Networking, DSA, Marvell, Atheros, MDIO,
I2C, SPI, UART, DDR, USB, Bluetooth, SD-Card, IPC, Multi-Thread, Socket,
Board Bring-Up, Porting, Device Tree, Zephyr, RTOS, Jenkins, TFTP,
Freescale, i.MX6, TI, TDA2x, NXP, S32R, IPQ4029, VVDN, PHAZR, Mirafra,
Randstad, Antal, Marlabs, OTA, Fail-Safe, Boot KPI, Performance Optimisation,
Embedded Systems, Firmware, Cache, PMIC, Clocks, PBL, XBL, GearVM
}
\end{center}

\end{document}
